\documentclass[10pt]{article}

\usepackage[letterpaper,margin=0.75in]{geometry}
\usepackage{amsmath,amssymb}
\usepackage{booktabs}
\usepackage{graphicx}
\usepackage[hidelinks]{hyperref}
\usepackage{xcolor}
\usepackage{tabularx}

% Put every image used by this report in report/figs/. Figure calls below
% therefore need only the filename, for example: density_grid.png.
\graphicspath{{figs/}}

\setlength{\parindent}{0pt}
\setlength{\parskip}{0.45em}
\setlength{\textfloatsep}{0.7em}
\setlength{\floatsep}{0.7em}
\setlength{\intextsep}{0.7em}

% A visible writing area. Replace each entire \answerbox call with prose when
% drafting the final report.
\newcommand{\answerbox}[2]{%
  \noindent\fcolorbox{gray!55}{gray!5}{%
    \begin{minipage}{0.965\linewidth}
      \textcolor{black!90}{\textit{#1}}
      \par\vspace{#2}
    \end{minipage}%
  }\par
}

% Render images side by side
\newcommand{\reportfigurepair}[4]{%
  \begin{figure}[tb]
  \centering
  \begin{minipage}{0.48\linewidth}
  \centering
  \IfFileExists{figs/#1}{%
    \includegraphics[width=\linewidth,height=0.27\textheight,keepaspectratio]{#1}%
  }{%
    \fbox{\parbox[c][0.16\textheight][c]{0.92\linewidth}{%
      \centering\textcolor{gray!75}{Figure placeholder\\[0.4em]
      Put \texttt{\detokenize{#1}} in \texttt{figs/}}}%
    }
  }
  \end{minipage}%
  \hfill
  \begin{minipage}{0.48\linewidth}
  \centering
  \IfFileExists{figs/#2}{%
    \includegraphics[width=\linewidth,height=0.27\textheight,keepaspectratio]{#2}%
  }{%
    \fbox{\parbox[c][0.16\textheight][c]{0.92\linewidth}{%
      \centering\textcolor{gray!75}{Figure placeholder\\[0.4em]
      Put \texttt{\detokenize{#2}} in \texttt{figs/}}}%
    }
  }
  \end{minipage}
  \caption{#3}
  \label{#4}
  \end{figure}
}

% Includes an image when it exists and otherwise leaves a compile-safe box.
% Usage: \reportfigure{filename.png}{Caption text.}{fig:label}
\newcommand{\reportfigure}[3]{%
  \begin{figure}[tb]
    \centering
    \IfFileExists{figs/#1}{%
      \includegraphics[width=0.88\linewidth,height=0.29\textheight,keepaspectratio]{#1}%
    }{%
      \fbox{\parbox[c][0.16\textheight][c]{0.82\linewidth}{%
        \centering\textcolor{black!90}{Figure placeholder\\[0.4em]
        Put \texttt{\detokenize{#1}} in \texttt{figs/}}}%
    }
    }
    \caption{#2}
    \label{#3}
  \end{figure}
}

\title{HW \#1: Mixture Density Networks}
\author{Nicole Lacey \and Christopher Lawton}
\date{Probabilistic Deep Learning --- Fall 2026}

\begin{document}
\maketitle
\vspace{-2.2em}


\section{Model and implementation}

For an input $\mathbf{x}$ and target $\mathbf{y}$, our mixture density network
(MDN) models the conditional density as
\begin{equation}
  p(\mathbf{y}\mid\mathbf{x}) =
  \sum_{k=1}^{K}\pi_k(\mathbf{x})
  \mathcal{N}\!\left(\mathbf{y};\boldsymbol{\mu}_k(\mathbf{x}),
  \boldsymbol{\Sigma}_k(\mathbf{x})\right),
  \qquad \sum_{k=1}^{K}\pi_k(\mathbf{x})=1.
\end{equation}
The network is trained by minimizing the mean negative log likelihood (NLL).
To implement this, we created an NN module that outputs the mixture weights (through a softmax), $\mu$, and $\Sigma$ (through a softplus, plus an additional term to ensure a minimum scale of $\epsilon>0$) for each mixture component. We optimize using the log-sum-exp expansion. Our hyperparameters include the number of mixture components $K$ (the number of Gaussians in the mixture), the rank of the covariance matrices (4), learning rate (1e-3), batch size (64), and number of epochs (2000). 

\section{Synthetic two-dimensional example}
The data we used for our 2D synthetic example was generated by sampling 4000 points from two different Gaussian distributions with different means and covariances. Then, we added Gaussian noise to the output. This resulted in a multimodal distribution for the output given the input, see Figure~\ref{fig:synthetic-data}. This resulting trained model should show a single Guassian near the intersection of the distributions, and two separate Gaussians that are increasingly far apart as the input moves away from the intersection (if the MDN learns well). Figure \ref{fig:density-grid-2d} does indeed show this behavior. The MDN adequately finds each mode with respect to the data.

Additionally, we tried a range of values for $K$, see Figure~\ref{fig:synthetic-comparison}, and found that there was no significant improvement for any $K$ greater than 2. This is expected because the data comes from a simple two-modal distribution.

  
We also compare these $K$ models to a deterministic regressor. This regressor performed similarly to the $K=1$ MDN. NLL and RMSE are reported for all the aforementioned models in Table~\ref{tab:synthetic-results}, when relevant. The test RMSE is lowest for $K=5$, at 0.9054, but $K=2$ is very close at 0.9056. 

We observe that training loss (NLL or MSE) over epochs decreases for all models, and that MDN models with $K>1$ outperform the deterministic and single-Gaussian models. Models with $K \in {2,3,4,5}$ perform similarly, with no clear best choice. 




\reportfigure{synthetic_raw_data.png}{Raw training observations for the
synthetic task, showing the multiple output modes associated with the input.}
{fig:synthetic-data}

\reportfigure{density_grid_2d.png}{Learned conditional density for the synthetic
2D task over a grid of input and output values. The density should reveal the
modes captured by the MDN and how their locations, weights, and widths vary with
the input.}{fig:density-grid-2d}



\reportfigure{synthetic_training.png}{Training objective and RMSE across epochs
for the deterministic baseline, single-Gaussian model, and MDNs with different
numbers of components.}
{fig:synthetic-comparison}

\begin{table}[tb]
  \centering
  \caption{Synthetic-data performance loaded directly from
  \texttt{csvs/synthetic\_summary.csv}. NLL is not defined for the
  deterministic model.}
  \label{tab:synthetic-results}
  \begin{tabular}{lccccc}
  \toprule
  Model & $K$ & Train NLL & Test NLL & Train RMSE & Test RMSE \\
  \midrule
  Deterministic & --- & N/A & N/A & 0.914328 & 0.920263 \\
  Single Gaussian & 1 & 0.832287 & 0.807886 & 0.921331 & 0.905627 \\
  MDN & 2 & -0.257680 & -0.263095 & 0.920119 & 0.907287 \\
  MDN & 3 & -0.263673 & -0.265676 & 0.919707 & 0.909541 \\
  MDN & 4 & -0.262717 & -0.265303 & 0.919435 & 0.910978 \\
  MDN & 5 & -0.262428 & -0.257620 & 0.920986 & 0.905491 \\
  \bottomrule
\end{tabular}

\end{table}

\newpage
\section{Higher-dimensional realistic data}

We use the Wind Turbine SCADA dataset available on Kaggle as the realistic data set.
This data set contains the features described in Table~\ref{tab:scada_features}.
To construct a trainable feature set we decided to drop the Theoretical Power Curve, as we are unsure how the manufacturer generated this feature.
We also replaced the raw Date/Time timestamp with cyclical (sine/cosine) encodings of hour-of-day and day-of-year, mapping each periodic time feature onto the unit circle via $\theta = 2\pi t / \text{period}$ (period = 24 for hour, using the fractional hour to account for the dataset's 10-minute sampling; period = 365 for day-of-year) so that adjacent times (e.g. 24:00 and 00:01) remain close in feature space rather than being separated by the full width of the raw numeric scale.
Wind speed and wind direction were kept as direct scalar inputs.

\begin{table}[htbp]
    \centering
    \caption{Features of the wind turbine SCADA dataset.}
    \label{tab:scada_features}
    \begin{tabularx}{\linewidth}{@{} l l X @{}}
        \toprule
        \textbf{Feature} & \textbf{Unit} & \textbf{Description} \\
        \midrule
        Date/Time & -- & Timestamp, recorded at 10-minute intervals \\
        LV ActivePower & kW & Power generated by the turbine at that moment \\
        Wind Speed & m/s & Wind speed at the hub height of the turbine \\
        Theoretical Power Curve & kWh & Theoretical power the turbine generates at a given wind speed, as specified by the manufacturer \\
        Wind Direction & $^\circ$ & Wind direction at the hub height of the turbine \\
        \bottomrule
    \end{tabularx}
\end{table}

Figure~\ref{fig:wind-features} shows each feature plotted against the target, LV ActivePower.
From a visual perspective, most features did not obviously display a multimodal relationship with respect to ActivePower except for Wind Speed.
In most cases the entire range of ActivePower can be observed for nearly all feature values, however, Wind Speed does display a multimodal relationship with the target: across the ramp region (roughly 4-13 m/s) and at or near maximum power output (~13 m/s and beyond).
Wind Direction shows a related but weaker pattern, power appears to switch between near-zero and near-max output across direction bands, in particular straddling 200°, but this on/off behavior is not by itself evidence of a genuinely multimodal conditional distribution.
Figure~\ref{fig:wind-conditional-histograms} shows the conditional distribution of ActivePower within narrow bands straddling several representative query points: $5\pm0.5$, $9\pm0.5$, and $13\pm0.5$ m/s for Wind Speed, $50\pm10^\circ$ and $200\pm10^\circ$ for Wind Direction.
Unlike the scatter in Figure~\ref{fig:wind-features}, this makes any separation in $P(\text{ActivePower}\mid x_0)$ explicit as distinct histogram peaks rather than something inferred from the shape of a scatter cloud.

\reportfigure{wind_turbine_features.png}{Wind-turbine active power against each
available numeric predictor, used to assess multimodality and
heteroscedasticity before fitting the models.}
{fig:wind-features}

\reportfigurepair{wind_turbine_conditional_histograms_speed.png}{wind_turbine_conditional_histograms_direction.png}
{Conditional histograms of ActivePower at fixed slices of wind speed (left) and wind
direction (right).}
{fig:wind-conditional-histograms}

\reportfigure{density_grid_wind_turbine.png}{Learned conditional density for the
wind-turbine task over a grid of inputs and active-power outputs. Use this view
to show heteroscedasticity, multiple operating regimes, and regions where the
model assigns diffuse or multimodal probability.}{fig:density-grid-wind}

\begin{table}[tb]
  \centering
  \caption{Wind-turbine performance loaded directly from
  \texttt{csvs/wind\_turbine\_summary.csv}. RMSE is reported in kW; NLL is
  evaluated on standardized targets and is not defined for the deterministic
  model.}
  \label{tab:wind-results}
  \begin{tabular}{lccccc}
  \toprule
  Model & $K$ & Train NLL & Test NLL & Train RMSE & Test RMSE \\
  \midrule
  Deterministic & --- & N/A & N/A & 383.646149 & 389.206085 \\
  Single Gaussian & 1 & -0.548006 & -0.558157 & 386.391083 & 391.403168 \\
  MDN & 2 & -1.100708 & -1.081633 & 391.576721 & 398.146759 \\
  MDN & 3 & -0.783956 & -0.764667 & 392.832428 & 399.270935 \\
  MDN & 4 & -1.279749 & -1.283014 & 391.055908 & 397.277557 \\
  MDN & 5 & -1.281521 & -1.287033 & 392.594330 & 398.612457 \\
  \bottomrule
\end{tabular}

\end{table}

\reportfigure{wind_turbine_training.png}{Training and validation/test behavior
for the realistic-data models, including the selected checkpoint or stopping
point.}{fig:wind-training}

\answerbox{Compare training and test behavior on the realistic dataset. Discuss
generalization gaps, predictive sharpness and calibration, regions of high
uncertainty, the effect of $K$, and how conclusions differ from the controlled
synthetic example. Reference Figures~\ref{fig:wind-features} and
\ref{fig:density-grid-wind}, Figure~\ref{fig:wind-training}, and
Table~\ref{tab:wind-results}.}{7\baselineskip}

\section{Conclusion}

\answerbox{Summarize when the MDN improves on deterministic and single-Gaussian
regression, identify the best-supported mixture count, state limitations, and
give one concrete next experiment.}{4\baselineskip}

% Uncomment this section if references are used.
% \begin{thebibliography}{9}
% \bibitem{bishop1994mdn}
% C. M. Bishop, ``Mixture Density Networks,'' 1994.
% \end{thebibliography}

\end{document}
